\section{课后练习}

\begin{frame}{练习一：实现\texttt{rstrip}}
接受一个\texttt{std::string}，返回删除末尾连续空白后的结果。
\begin{enumerate}
\item 空串应直接返回，不能在空串上调用\texttt{back()}或\texttt{pop\_back()}。
\item 可以按值接受参数，在局部副本上修改，调用者原字符串保持不变。
\item 使用\texttt{std::isspace}时，把字符转换为\texttt{unsigned char}后再传入。
\end{enumerate}
\medskip
\textbf{示例：}
\begin{itemize}
\item 输入\texttt{"hello world\ \ \ "}，返回\texttt{"hello world"}。
\item 输入只有空白时，返回空串。
\end{itemize}
\alert{先写出循环继续条件，再检查它是否总能保证\texttt{back()}有效。}
\end{frame}

\begin{frame}{练习一参考实现}
\begin{Code}*{cpp}\lstinputlisting{code/rstrip.cpp}\end{Code}
\small 按值传参是为了得到可修改副本；返回局部对象不需要手动释放。
\end{frame}

\begin{frame}{练习二：打印每个单词中的大写字母}
给定\texttt{std::vector<std::string>}，逐行输出每个字符串中的大写字母。
\begin{enumerate}
\item 外层遍历只读取单词，使用\texttt{const std::string \&}避免逐个复制。
\item 内层遍历当前字符串中的字符。
\item 使用\texttt{std::isupper}判断，并正确转换函数参数。
\end{enumerate}
\medskip
\textbf{示例数据：}\texttt{\{"Hello", "World", "C++"\}}

\textbf{预期输出：}三行分别为\texttt{H}、\texttt{W}、\texttt{C}。
\medskip
\alert{嵌套范围for可以直接表达“每个单词中的每个字符”。}
\end{frame}

\begin{frame}{练习二参考实现}
\begin{Code}*{cpp}\lstinputlisting{code/uppercase.cpp}\end{Code}
\small 这里的外层元素类型是\texttt{std::string}，内层元素类型是\texttt{char}；两种遍历使用不同的拷贝策略。
\end{frame}

\begin{frame}{迁移自查：这些写法真的等价吗？}
\small
\setlength{\tabcolsep}{4pt}
\renewcommand{\arraystretch}{1.3}
\begin{tabular}{p{0.34\textwidth}p{0.50\textwidth}}
\textbf{习惯或误写} & \textbf{更合适的C++表达／需要补充的检查}\\
\hline
\texttt{strlen(s.c\_str())} & \texttt{s.size()}\\
\texttt{strcmp(a.c\_str(), b.c\_str())} & \texttt{a == b}或相应关系运算\\
\texttt{while (!v.empty()) v.pop\_back();} & \texttt{v.clear()}\\
\texttt{reserve(n)}后用\texttt{v[0]} & \texttt{reserve}不创建元素；按需求使用\texttt{resize(n)}或追加元素\\
范围for中向同一容器追加 & 改为固定原范围或延迟修改\\
\texttt{using namespace std}遍布全局 & 明确写\texttt{std::}，或只引入需要的名字
\end{tabular}

\medskip\normalsize
\alert{更短的写法只有在边界、生命周期和语义都正确时才更好。}
\end{frame}

\begin{frame}{继续查阅}
\begin{itemize}
\item \href{https://en.cppreference.com/w/cpp/io}{\texttt{cppreference：输入输出库}}
\item \href{https://en.cppreference.com/w/cpp/string/basic_string}{\texttt{cppreference：std::string}}
\item \href{https://en.cppreference.com/w/cpp/container/vector}{\texttt{cppreference：std::vector}}
\item \href{https://en.cppreference.com/w/cpp/container}{\texttt{cppreference：容器库}}
\end{itemize}
\medskip
查阅成员函数时重点看四件事：
\begin{enumerate}
\item 参数与返回类型；
\item 前置条件与异常；
\item 复杂度；
\item 指针、引用和迭代器失效规则。
\end{enumerate}
\end{frame}
